\documentclass{practicaEPCC}

% --- PAQUETES DE IDIOMA Y CODIFICACION ---
\usepackage[utf8]{inputenc}
\usepackage[T1]{fontenc}
\usepackage[spanish]{babel}
% -----------------------------------------

\usepackage{hyperref}
\usepackage{amssymb}
\usepackage{amsmath}
\usepackage{minted}
\usepackage{multicol}
\usepackage{pgfplots}
\usepackage{graphicx}
\usepackage{float}
\usepackage{listings}
\usepackage{color}
\usepackage{adjustbox}
\usepackage{caption}
\usepackage{subcaption}
\usepackage{geometry}
\usepackage{enumitem}
\usepackage[strings]{underscore}

\pgfplotsset{compat=1.18}

\lstset{
 basicstyle=\ttfamily\footnotesize,
 keywordstyle=\color{blue},
 commentstyle=\color{gray},
 stringstyle=\color{red},
 numbers=left,
 numberstyle=\tiny,
 stepnumber=1,
 numbersep=5pt,
 frame=single,
 tabsize=2,
 showstringspaces=false,
 breaklines=true,
 captionpos=b
}

\newcommand{\imageplaceholder}[2]{
\begin{figure}[H]
\centering
\fbox{
\begin{minipage}[c][7cm][c]{0.9\textwidth}
\centering
\vspace{0.5cm}
\textbf{Espacio para imagen}\\[0.4cm]
#1\\[0.4cm]
\textit{Reemplazar este recuadro por:}\\
\texttt{\string\includegraphics[width=0.95\string\textwidth]\{Imagenes/#2\}}
\vspace{0.5cm}
\end{minipage}
}
\caption{#1}
\end{figure}
}

\school{Escuela Profesional de Ciencia de la Computacion}
\course{Big Data}
\professor{}
\title{Laboratorio: Instalacion y demostracion de Apache Kafka sobre instancias EC2 en AWS}
\author{Michael Jarnie Ticona Larico}
\date{\today}

\begin{document}

\maketitle

\section{Instalacion de Apache Kafka en AWS EC2}

\subsection{Objetivo}

El objetivo del laboratorio fue instalar Apache Kafka sobre instancias EC2 en AWS y demostrar los conceptos fundamentales del modelo de mensajeria distribuida de Kafka: \textit{Producer}, \textit{Consumer}, \textit{Topic}, \textit{Partition}, \textit{Offset} y \textit{Consumer Group}.

Para automatizar el proceso se implemento un script en Python con \texttt{boto3}, encargado de crear las instancias EC2, inyectar scripts de inicializacion mediante \textit{user-data} y configurar automaticamente el nodo master y los nodos trabajadores.

\subsection{Arquitectura implementada}

La arquitectura utilizada considera una instancia inicial desde la cual se ejecuta el script de aprovisionamiento. Esta instancia actua como controlador u orquestador, pero no forma parte del cluster Kafka. Desde ella se crean nuevas instancias EC2 para el laboratorio.

Al ejecutar el comando:

\begin{lstlisting}[language=bash]
python3 src/levantar_kafka.py --start_nodes 2
\end{lstlisting}

se crean tres instancias EC2 nuevas:

\begin{itemize}
\item \texttt{Kafka-Master-<run\_id>}: nodo principal donde se instala y ejecuta el broker de Kafka.
\item \texttt{Kafka-Worker-1-<run\_id>}: nodo trabajador que actua como cliente Kafka.
\item \texttt{Kafka-Worker-2-<run\_id>}: segundo nodo trabajador que actua como cliente Kafka.
\end{itemize}

Cada instancia recibe un nombre unico mediante la etiqueta \texttt{Name}. Ademas, se agregan etiquetas como \texttt{Proyecto}, \texttt{Rol} y \texttt{ClusterRun}, permitiendo identificar facilmente los recursos creados en la consola de AWS.

\imageplaceholder{Instancias EC2 creadas para el laboratorio Kafka: master y workers con nombres unicos.}{kafka_instancias_ec2.png}

\subsection{Configuracion previa}

La configuracion de AWS se centralizo en el archivo \texttt{src/config.py}. Este archivo contiene los identificadores necesarios para crear las instancias EC2:

\begin{lstlisting}[language=Python]
REGION = 'us-east-1'
AMI_ID = 'ami-03f4fd1e8233bd64d'
KEY_NAME = 'cluster'
SECURITY_GROUP_ID = 'sg-00c82fc157b6a0478'
SUBNET_ID = 'subnet-0346fd19f61aafdcd'
TIPO_INSTANCIA = 't3.small'
USUARIO_SSH = 'ec2-user'
\end{lstlisting}

El grupo de seguridad debe permitir la comunicacion por SSH y por los puertos utilizados por Kafka:

\begin{itemize}
\item Puerto \texttt{22}: conexion SSH.
\item Puerto \texttt{9092}: conexion de clientes al broker Kafka.
\item Puerto \texttt{9093}: comunicacion interna del controlador Kafka en modo KRaft.
\end{itemize}

\subsection{Script de aprovisionamiento}

La creacion de las instancias se realizo mediante el script \texttt{src/levantar\_kafka.py}. Este programa utiliza \texttt{boto3} para comunicarse con AWS EC2 y crear los nodos necesarios.

Las acciones principales del script son:

\begin{enumerate}
\item Leer los scripts \texttt{user\_data\_master.sh} y \texttt{user\_data\_worker.sh}.
\item Crear una instancia EC2 para el master Kafka.
\item Obtener la IP privada del master.
\item Crear los workers inyectando la IP privada del broker.
\item Asignar nombres diferentes a cada nodo creado.
\item Mostrar los comandos necesarios para conectarse y ejecutar la demostracion.
\end{enumerate}

\begin{lstlisting}[language=bash]
python3 src/levantar_kafka.py --start_nodes 2
\end{lstlisting}

\imageplaceholder{Ejecucion del script de creacion de instancias Kafka desde Python.}{kafka_start_nodes.png}

\subsection{Instalacion automatica en el nodo master}

El nodo master recibe el script \texttt{scripts/user\_data\_master.sh}. Este script se ejecuta automaticamente al iniciar la instancia EC2 y realiza las siguientes tareas:

\begin{itemize}
\item Actualiza el sistema operativo.
\item Instala Java 11, \texttt{wget}, \texttt{tar} y \texttt{nmap-ncat}.
\item Descarga Apache Kafka 3.7.1.
\item Configura Kafka en modo KRaft, sin depender de ZooKeeper.
\item Crea el servicio \texttt{kafka.service} en \texttt{systemd}.
\item Inicia Kafka automaticamente.
\item Crea el script \texttt{/home/ec2-user/kafka\_demo.sh} para la demostracion.
\end{itemize}

\begin{lstlisting}[language=bash]
sudo systemctl status kafka
\end{lstlisting}

La salida obtenida muestra que Kafka quedo activo y escuchando conexiones en el puerto \texttt{9092}.

\imageplaceholder{Estado del servicio Kafka activo en el nodo master.}{kafka_systemctl_status.png}

\subsection{Instalacion automatica en los workers}

Los nodos trabajadores reciben el script \texttt{scripts/user\_data\_worker.sh}. A diferencia del master, estos nodos no ejecutan el broker Kafka, sino que instalan las herramientas cliente necesarias para conectarse al broker.

Cada worker queda configurado con la variable de entorno:

\begin{lstlisting}[language=bash]
KAFKA_BROKER=<IP_PRIVADA_DEL_MASTER>:9092
\end{lstlisting}

Ademas, se crea el script:

\begin{lstlisting}[language=bash]
/home/ec2-user/kafka_client_demo.sh
\end{lstlisting}

Este script permite ejecutar consumidores desde los workers para demostrar el uso de \textit{Consumer Groups}.

\imageplaceholder{Worker Kafka conectado al broker mediante la variable KAFKA\_BROKER.}{kafka_worker_cliente.png}

\section{Demostracion de conceptos de Kafka}

\subsection{Verificacion del broker}

Una vez finalizado \textit{cloud-init}, se verifico que Kafka estuviera ejecutandose correctamente con el siguiente comando:

\begin{lstlisting}[language=bash]
sudo systemctl status kafka
\end{lstlisting}

La salida mostro que el servicio estaba en estado \texttt{active (running)} y que Kafka habia iniciado correctamente:

\begin{lstlisting}[language={}]
Kafka Server started
Awaiting socket connections on 0.0.0.0:9092
Kafka version: 3.7.1
\end{lstlisting}

Esto confirma que el broker quedo listo para recibir productores y consumidores.

\subsection{Ejecucion de la demostracion}

La demostracion completa se ejecuto desde el master mediante:

\begin{lstlisting}[language=bash]
/home/ec2-user/kafka_demo.sh
\end{lstlisting}

Este script automatiza la creacion del topic, la publicacion de mensajes, el consumo de registros y la consulta de offsets por grupo de consumidores.

\imageplaceholder{Ejecucion completa del script kafka\_demo.sh en el master.}{kafka_demo_completa.png}

\subsection{Topic}

Un \textit{Topic} en Kafka es una categoria o canal logico donde se almacenan mensajes. En el laboratorio se creo el topic \texttt{onpe-votos-demo}.

\begin{lstlisting}[language=bash]
kafka-topics.sh --bootstrap-server "$KAFKA_BROKER" \
  --create --if-not-exists \
  --topic "onpe-votos-demo" \
  --partitions 3 \
  --replication-factor 1
\end{lstlisting}

Este comando crea el canal donde los productores escriben mensajes y desde donde los consumidores los leen.

\imageplaceholder{Creacion del topic onpe-votos-demo en Kafka.}{kafka_topic_creado.png}

\subsection{Partition}

Las \textit{Partitions} dividen un topic en segmentos independientes, permitiendo paralelismo y distribucion de carga. En este caso, el topic se creo con tres particiones.

Para verificarlo se uso:

\begin{lstlisting}[language=bash]
kafka-topics.sh --bootstrap-server "$KAFKA_BROKER" \
  --describe \
  --topic "onpe-votos-demo"
\end{lstlisting}

La salida muestra las particiones \texttt{0}, \texttt{1} y \texttt{2}. Cada mensaje publicado se almacena en una de estas particiones.

\imageplaceholder{Descripcion del topic mostrando sus tres particiones.}{kafka_particiones.png}

\subsection{Producer}

Un \textit{Producer} es el componente encargado de publicar mensajes en un topic. En el laboratorio se utilizo \texttt{kafka-console-producer.sh} para enviar mensajes simulando resultados de mesas electorales.

\begin{lstlisting}[language=bash]
cat <<MESSAGES | kafka-console-producer.sh \
  --bootstrap-server "$KAFKA_BROKER" \
  --topic "onpe-votos-demo" \
  --property parse.key=true \
  --property key.separator=:
LIMA:mesa=001 votos=120
AREQUIPA:mesa=002 votos=95
CUSCO:mesa=003 votos=110
LIMA:mesa=004 votos=130
AREQUIPA:mesa=005 votos=102
CUSCO:mesa=006 votos=99
MESSAGES
\end{lstlisting}

Se usaron llaves como \texttt{LIMA}, \texttt{AREQUIPA} y \texttt{CUSCO} para que Kafka distribuya los mensajes entre particiones de acuerdo con la llave.

\imageplaceholder{Producer enviando mensajes al topic onpe-votos-demo.}{kafka_producer.png}

\subsection{Consumer}

Un \textit{Consumer} lee los mensajes almacenados en un topic. Para consumir los mensajes desde el inicio se utilizo:

\begin{lstlisting}[language=bash]
kafka-console-consumer.sh \
  --bootstrap-server "$KAFKA_BROKER" \
  --topic "onpe-votos-demo" \
  --from-beginning \
  --timeout-ms 10000 \
  --property print.partition=true \
  --property print.offset=true \
  --property print.key=true \
  --property key.separator=" | "
\end{lstlisting}

La salida permite observar el contenido de los mensajes junto con su particion y offset.

\imageplaceholder{Consumer leyendo mensajes desde el topic Kafka.}{kafka_consumer.png}

\subsection{Offset}

El \textit{Offset} representa la posicion de un mensaje dentro de una particion. Kafka asigna un offset incremental por cada mensaje almacenado.

En la demostracion, el consumidor imprime el offset usando:

\begin{lstlisting}[language=bash]
--property print.offset=true
\end{lstlisting}

Por ejemplo, una salida puede mostrar mensajes con el siguiente formato:

\begin{lstlisting}[language={}]
Partition:0 | Offset:0 | LIMA | mesa=001 votos=120
Partition:1 | Offset:0 | AREQUIPA | mesa=002 votos=95
Partition:2 | Offset:0 | CUSCO | mesa=003 votos=110
\end{lstlisting}

Esto demuestra que cada particion mantiene su propia secuencia de offsets.

\imageplaceholder{Salida del consumer mostrando particion y offset de cada mensaje.}{kafka_offsets.png}

\subsection{Consumer Group}

Un \textit{Consumer Group} permite que varios consumidores trabajen de manera coordinada leyendo un mismo topic. Kafka distribuye las particiones entre los consumidores del grupo.

En el laboratorio se uso el grupo \texttt{grupo-consulta-onpe}:

\begin{lstlisting}[language=bash]
kafka-console-consumer.sh \
  --bootstrap-server "$KAFKA_BROKER" \
  --topic "onpe-votos-demo" \
  --group "grupo-consulta-onpe" \
  --from-beginning \
  --max-messages 6
\end{lstlisting}

Luego se consulto el estado del grupo y sus offsets confirmados:

\begin{lstlisting}[language=bash]
kafka-consumer-groups.sh \
  --bootstrap-server "$KAFKA_BROKER" \
  --describe \
  --group "grupo-consulta-onpe"
\end{lstlisting}

La salida muestra el \texttt{CURRENT-OFFSET}, el \texttt{LOG-END-OFFSET} y el \texttt{LAG}, lo que permite verificar cuantos mensajes fueron consumidos por el grupo.

\imageplaceholder{Descripcion del consumer group mostrando offsets y lag.}{kafka_consumer_group.png}

\section{Prueba desde los nodos workers}

Para demostrar que los workers pueden actuar como clientes Kafka, se accedio por SSH a un worker y se ejecuto:

\begin{lstlisting}[language=bash]
/home/ec2-user/kafka_client_demo.sh onpe-votos-demo grupo-workers-onpe
\end{lstlisting}

Si se ejecuta el mismo comando en dos workers diferentes con el mismo grupo, Kafka reparte las particiones entre ambos consumidores. Esto evidencia la coordinacion propia de los \textit{Consumer Groups}.

\imageplaceholder{Consumers ejecutados desde workers distintos usando el mismo Consumer Group.}{kafka_workers_consumer_group.png}

\section{Eliminacion de recursos}

Para evitar costos innecesarios en AWS, al finalizar el laboratorio se eliminaron las instancias creadas mediante:

\begin{lstlisting}[language=bash]
python3 src/levantar_kafka.py --delete
\end{lstlisting}

Este comando busca las instancias con la etiqueta \texttt{Proyecto=ONPE-Kafka} y las termina automaticamente.

\imageplaceholder{Eliminacion de las instancias EC2 creadas para Kafka.}{kafka_delete.png}

\section{Analisis}

Apache Kafka funciona como una plataforma distribuida de mensajeria basada en logs. A diferencia de una cola tradicional, los mensajes no desaparecen inmediatamente al ser consumidos; permanecen almacenados en particiones y los consumidores avanzan mediante offsets.

El uso de particiones permite paralelizar la lectura y escritura de datos. En este laboratorio, el topic \texttt{onpe-votos-demo} fue dividido en tres particiones, lo que permitio observar como Kafka organiza los mensajes internamente.

El concepto de \textit{Consumer Group} es especialmente importante en escenarios de Big Data, ya que permite distribuir el procesamiento entre multiples consumidores. Si varios workers pertenecen al mismo grupo, Kafka asigna particiones distintas a cada uno, evitando que todos procesen exactamente los mismos mensajes.

La automatizacion con \texttt{boto3} permitio crear una infraestructura reproducible sobre AWS EC2. El master quedo configurado como broker Kafka, mientras que los workers quedaron preparados como clientes para realizar pruebas de consumo.

\section{Conclusiones}

Se logro instalar Apache Kafka sobre instancias EC2 en AWS utilizando scripts automatizados de aprovisionamiento e inicializacion.

La instancia local o inicial actuo como controlador, creando tres nodos EC2 nuevos: un master Kafka y dos workers cliente. Cada nodo fue creado con un nombre unico para facilitar su identificacion en AWS.

La demostracion permitio comprobar los conceptos fundamentales de Kafka: creacion de topics, division en partitions, envio de mensajes mediante producers, lectura mediante consumers, seguimiento de offsets y coordinacion de consumidores mediante consumer groups.

Finalmente, el laboratorio mostro que Kafka es una herramienta adecuada para construir sistemas distribuidos de ingestion y procesamiento de eventos, especialmente cuando se requiere escalabilidad, tolerancia a fallos y consumo paralelo de datos.

\end{document}
